\newpage

\stepcounter{section}
\section*{Lecture 14: Products}
Let $X,Y$ be schemes. In the same way that we can take the product of two sets or groups we want to define a product scheme $X\times Y$. This could mean taking the set of points $X\times Y$ and putting a scheme structure on it, i.e. topologising the set and providing a sheaf of rings. However this turns out to not be ``the right thing'' to do.
To see why, we must first pin down what we mean by ``the right thing'' in this context.

To determine the correct definition we must specify the properties we want the product to satisfy. Our intuition tells us that we should have maps $\pi_X\colon X\times Y \to X$ and $\pi_Y\colon X\times Y \to Y$ such that for any scheme $Z$ we have
\[
\text{Maps}(Z, X\times Y)\to \text{Maps}(Z,X) \times \text{Maps}(Z,Y),
\]
\[
\varphi\mapsto (\pi_X\circ \varphi, \pi_Y\circ \varphi),
\]
and this assignment is a bijection. This leads us to define the product as the triple $(X\times Y, \pi_X, \pi_Y)$ that satisfies the following commutative diagram.
\[\begin{tikzcd}
	&&& X \\
	Z && {X\times Y} \\
	&&& Y
	\arrow[curve={height=-6pt}, from=2-1, to=1-4]
	\arrow["{\exists!}", dashed, from=2-1, to=2-3]
	\arrow[curve={height=6pt}, from=2-1, to=3-4]
	\arrow["{\pi_X}"', from=2-3, to=1-4]
	\arrow["{\pi_Y}", from=2-3, to=3-4]
\end{tikzcd}\]
We then ask two questions,
\begin{enumerate}
    \item Does such an object exist?
    \item Is such an object unique?
\end{enumerate}
The answer to the second one is typically false, as any object satisfying the diagram can be relabelled, so we instead ask if such an object is unique up to unique isomorphism. That is, if $\square$ and $\bigcirc$ are two such specified objects there exists a unique isomorphism $\square\to\bigcirc$ compatible with the projection maps.

Suppose $\square$ and $\bigcirc$ are two objects that fit the commutative diagram. Then we can let $Z = \square$ to obtain that the identity 
id$_\square$ is the unique map which makes the diagram commute. We then use the fact that $\bigcirc$ is also a product to obtain the following diagram,
\[\begin{tikzcd}
	&&&&& X \\
	\square && \bigcirc && \square \\
	&&&&& Y
	\arrow[curve={height=-18pt}, from=2-1, to=1-6]
	\arrow["{\exists! f}", dashed, from=2-1, to=2-3]
	\arrow[curve={height=18pt}, from=2-1, to=3-6]
	\arrow[curve={height=-6pt}, from=2-3, to=1-6]
	\arrow["{\exists! g}", dashed, from=2-3, to=2-5]
	\arrow[curve={height=6pt}, from=2-3, to=3-6]
	\arrow[from=2-5, to=1-6]
	\arrow[from=2-5, to=3-6]
\end{tikzcd}\]
so $g\circ f$ is also a map which works. By uniqueness of the map we obtain Id$_\square = g\circ f$, and similarly by swapping $\square$ and $\bigcirc$ we obtain Id$_\bigcirc = f\circ g$. Thus $f$ is the unique isomorphism from $\square$ to $\bigcirc$ as desired.

We have shown that if a product exists then it is unique up to unique isomorphism, so we denote the product by $X\times Y$. We now consider the question of existence.

\begin{example}
(Sets)
Let $X,Y$ be sets, we want to define the product $(X\times Y, \pi_X, \pi_Y)$. We claim that the Cartesian product
\[
X\times Y \colonequals \{ (x,y) \mid x\in X, y\in Y\},
\]
with $\pi_X,\pi_Y$ the projection maps to the first and second coordinate respectively, satisfies the specification. For any test set $T$ with maps $f,g$ we obtain the diagram,
\[\begin{tikzcd}
	&&& X \\
	T && {X\times Y} \\
	&&& Y
	\arrow["f", curve={height=-30pt}, from=2-1, to=1-4]
	\arrow["{t\mapsto(f(t),g(t))}", from=2-1, to=2-3]
	\arrow["g"', curve={height=30pt}, from=2-1, to=3-4]
	\arrow[from=2-3, to=1-4]
	\arrow[from=2-3, to=3-4]
\end{tikzcd}\]
where the uniqueness of the middle arrow is because for any $\varphi\colon T \to X\times Y$ satisfying the diagram we must have $\pi_X\circ \varphi = f$ and $\pi_Y\circ \varphi = g$, so $\varphi(t) = (f(t),g(t))$.
\end{example}

\begin{example}\hfill
\begin{itemize}
    \item (Topological spaces) Let $X,Y$ be topological spaces. Then $X\times Y$ is the Cartesian product with the product topology.
    \item (Groups) Let $G,H$ be groups. Then $G\times H$ is the Cartesian product with component-wise multiplication.
\end{itemize}
\end{example}

In all the examples we have done thus far the product is the Cartesian product with some additional structure imposed on it.
Surprisingly, this is not the case with schemes.
Here is the first example.

\begin{example}
    Let $X=Y=\Spec \Z$. Then the data of a test scheme $Z$ with maps $f\colon Z\to X$, $g\colon Z\to Y$ is equivalent to ring maps $\varphi\colon \Z\to \mathcal{O}_Z(Z)$, $\rho\colon \Z\to \mathcal{O}_Z(Z)$. However there is a unique ring map from $\Z$ to any ring, so $\rho=\varphi$ and the data of the map from $Z$ to $X\times Y$ is equivalent to the ring map $\varphi$. Thus we obtain that $X\times Y = \Spec \Z$, which satisfies the commutative diagram
    \[\begin{tikzcd}
	&&& {\Spec \Z} \\
	Z && {\Spec \Z} \\
	&&& {\Spec \Z}
	\arrow["f", from=2-1, to=1-4]
	\arrow[dashed, from=2-1, to=2-3]
	\arrow["g"', from=2-1, to=3-4]
	\arrow[from=2-3, to=1-4]
	\arrow[from=2-3, to=3-4]
\end{tikzcd}\]
\end{example}

\begin{example}
    Let $X = \Spec (\Z/2\Z)$ and $Y = \Spec (\Z/3\Z)$. Then $X\times Y = \emptyset$ is the empty scheme. This is because a scheme $Z$ with maps $f\colon Z\to X$, $g\colon Z\to Y$ is equivalent to ring maps $\varphi\colon \F_2\to \mathcal{O}_Z(Z)$, $\rho\colon \F_3\to \mathcal{O}_Z(Z)$. However we observe that
    \[
    1_Z = 3\rho(1) - 2\varphi(1) = \rho(3) - \varphi(2) = 0_Z,
    \]
    so $\mathcal{O}_Z(Z)$ must be the zero ring, and therefore \(Z\) must be the empty scheme.
    In other words, the only scheme that maps to \(X\) and \(Y\) is the empty scheme.
    As a result, the product \(X \times Y\) must also be the empty scheme, so we obtain the diagram
    \[\begin{tikzcd}
	&&& {\Spec \F_2} \\
	\emptyset && \emptyset \\
	&&& {\Spec \F_3}
	\arrow["f", from=2-1, to=1-4]
	\arrow[dashed, from=2-1, to=2-3]
	\arrow["g"', from=2-1, to=3-4]
	\arrow[from=2-3, to=1-4]
	\arrow[from=2-3, to=3-4]
\end{tikzcd}\]
\end{example}

\begin{example}
    Let $X = \Spec \Z[x] $ and $Y = \Spec \Z[y]$ be two copies of $\A^1_\Z $. Then $X\times Y = \Spec \Z[x,y] = \A_\Z^2$. To see why let $T$ be a test scheme with maps $ T\to X$, $T\to Y$, then this data is equivalent to ring maps $\varphi\colon \Z[x] \to \mathcal{O}_T(T)$, $\rho\colon \Z[y]\to \mathcal{O}_T(T)$. However the data of $(\varphi,\rho)$ is the same as choosing $x\mapsto f\in \mathcal{O}_T(T)$ and $y\mapsto g\in \mathcal{O}_T(T)$, which uniquely specifies a map $\Z[x,y] \to\mathcal{O}_T(T)$. All this information is then captured by the following diagram.
    \[\begin{tikzcd}
	&&& {\Spec \Z[x]} \\
	T && {\Spec \Z[x,y]} \\
	&&& {\Spec \Z[y]}
	\arrow["{x\mapsto f}", curve={height=-24pt}, from=2-1, to=1-4]
	\arrow["{(x,y)\mapsto(f,g)}", dashed, from=2-1, to=2-3]
	\arrow["{y\mapsto g}"', curve={height=24pt}, from=2-1, to=3-4]
	\arrow[from=2-3, to=1-4]
	\arrow[from=2-3, to=3-4]
\end{tikzcd}\]

By the same argument we also have for $\A^n_\Z = \Spec \Z[x_1,\dots, x_n]$ and $\A^m_\Z = \Spec \Z[y_1,\dots, x_m]$ that $\A^n_\Z\times \A^m_\Z = \A^{n+m}_\Z$.
\end{example}

We now note that the set of prime ideals of \(\Z[x,y]\) is not the product of the set of prime ideals of \(\Z[x]\) and \(\Z[y]\).
That is, we do not have the equality of sets
\[
\Spec \Z[x,y] = \Spec \Z[x] \times \Spec \Z[y].
\]

All the examples above contradict our intuition that the points of a product should be the product of the points of the parts.
Nevertheless, the intuition is correct if we use a different notion of points than the literal points of the underlying topological space.

We recall the two notions of points of a scheme $X$:
\begin{enumerate}
    \item The actual points of $X$ as a topological space,
    \item Given a field (or even a ring) \(k\), the $k$-points of \(X\), namely
      \[ \operatorname{Maps}(\Spec k, X).\]
\end{enumerate}
We often think of the first type of points as an implementation detail, defined purely to get a topology that admits a good sheaf of rings.
The second notion of points is more fundamental.
For example, suppose \(X\) is the affine scheme \(\Spec \Z[x_1,\dots,x_n]/(f_1,\dots,f_m)\).
Then the set of \(k\)-points of \(X\) is precisely the set of \(k\)-valued solutions to the system of equations
\[ f_1(x_1,\dots,x_n) = 0, \dots, f_m(x_1,\dots,x_n) = 0.\]
The primary concern of algebraic geometry is understanding solutions of polynomial equations!

The definition of the product scheme \(X \times Y\) behaves exactly as intended for this second, more fundamental notion of a point.
By definition, we have
\[ \operatorname{Maps}(\Spec k, X \times Y) = \operatorname{Maps}(\Spec k, X) \times \operatorname{Maps}(\Spec k,Y). \]
For any \(k\), the \(k\)-valued points of \(X \times Y\) are precisely the product of the \(k\)-valued points of \(X\) and \(Y\).

We now turn our attention to the general question: given \(X\) and \(Y\), does \(X \times Y\) exist?
Let us take the simpler case where both \(X\) and \(Y\) are affine, say $\Spec A$ and $Y =\Spec B$.
Then asking for \(X \times Y\) is asking whether there is a scheme $\square$ which satisfies the universal property in the following diagram: 
\[\begin{tikzcd}
	&&& {\Spec A} \\
	T && \square \\
	&&& {\Spec B}
	\arrow["{A\to \mathcal{O}_T(T)}", curve={height=-0pt}, from=2-1, to=1-4]
	\arrow[dashed, from=2-1, to=2-3]
	\arrow["{B\to \mathcal{O}_T(T)}"', curve={height=0pt}, from=2-1, to=3-4]
	\arrow[from=2-3, to=1-4]
	\arrow[from=2-3, to=3-4]
      \end{tikzcd}.\]
    We can try to find an affine scheme \(\square\) that satisfies the property.
    Trying to find an affine \(\square\) is equivalent to trying to find a ring $R$ with maps $A\to R$ and $B\to R$ satisfying the following universal property:
\[\begin{tikzcd}
	&&& A \\
	S && R \\
	&&& B
	\arrow[from=1-4, to=2-1]
	\arrow[from=1-4, to=2-3]
	\arrow[dashed, from=2-3, to=2-1]
	\arrow[from=3-4, to=2-1]
	\arrow[from=3-4, to=2-3]
\end{tikzcd}\]

Such a ring does exist and we will now construct it.
Consider the set 
\[
G \colonequals \{ a\otimes b \mid a\in A, b\in B\}.
\]
Let \(A \otimes B\) be the free abelian group generated by \(G\) modulo the following relations:
\[
  \begin{split}
    a_1 \otimes b + a_2\otimes b - (a_1 +a_2)\otimes b, \text{ for all } a_1, a_2 \in A \text{ and } b \in B,\\
    a\otimes b_1 + a\otimes b_2 - a \otimes (b_1 + b_2), \text{ for all } a \in A \text{ and } b_1,b_2 \in B.
  \end{split}
\]
We make \(A \otimes B\) a ring by setting
\[ (a_1 \otimes b_1) \times (a_2 \otimes b_2) = a_1a_2 \otimes b_1b_2,\]
and extending by linearity.
We need to check that this is well-defined modulo the relations, and satisfies the ring axioms, but we skip those checks.
We define maps $A \to A\otimes B$ and $B \to A\otimes B$ by $a\mapsto a\otimes 1$ and $b\mapsto 1\otimes b$ respectively.
Again, it is easy to check that these are ring homomorphisms.

Given maps $f\colon A \to S$ and $g\colon B \to S$, there is a unique map \(h \colon A \otimes B \to S\) such that the composite \(A \to A \otimes B \to S\) is \(f\) and \(B \to A \otimes B \to S\) is \(g\).
This map is given by
\[ h(a \otimes b) = f(a) g(b).\]
It is easy to check that this is a well-defined homomorphism.

We conclude that $\Spec A \times \Spec B$ exists and it equal to $\Spec(A\otimes B)$. As an exercise show that $\F_2\otimes \F_3 = 0$, which agrees with what we expected from before.
%%% Local Variables:
%%% mode: LaTeX
%%% TeX-master: "main"
%%% End:
